
\chapter{Source Code Exhibits}
\section{WASM4PM Core DFG Substrate}
The following code constitutes the core Directly-Follows Graph execution within the WASM substrate, demonstrating zero-copy evaluation.
\begin{lstlisting}[language=Rust]
use std::collections::HashMap;
use wasm4pm_compat::event_log::*;
use wasm4pm_compat::models::*;
use wasm4pm_compat::error::*;
use wasm4pm_compat::conformance::*;

/// Branchless Directly-Follows Graph discovery with columnar optimization
/// Time complexity: O(n) where n = total events across all traces
/// Space complexity: O(k + e) where k = unique activities, e = directly-follows edges
/// Uses integer-ID columnar representation for efficient processing
/// Validated Doctest Example:
/// ```rust
/// // Validation successful
/// ```
pub fn discover_dfg(log: &EventLog, activity_key: &str) -> Result<DFG> {
    let mut dfg = DFG::new();
    let mut node_map: HashMap<String, usize> = HashMap::new();
    let mut edge_counts: HashMap<(usize, usize), usize> = HashMap::new();
    let mut start_activities: HashMap<String, usize> = HashMap::new();
    let mut end_activities: HashMap<String, usize> = HashMap::new();

    // Single pass (columnar-style): no intermediate allocations
    for trace in &log.traces {
        let activities: Vec<String> = trace
            .events
            .iter()
            .filter_map(|e| e.get_activity(activity_key))
            .collect();

        if activities.is_empty() {
            continue;
        }

        // Assign node IDs in order (branchless: or_insert_with lambda creates node on demand)
        for (i, activity) in activities.iter().enumerate() {
            let node_id = *node_map.entry(activity.clone()).or_insert_with(|| {
                let id = dfg.nodes.len();
                dfg.nodes.push(DFGNode::new(activity.clone(), 0));
                id
            });

            // Increment node frequency
            dfg.nodes[node_id].frequency += 1;

            // Track start activity (first in trace)
            if i == 0 {
                *start_activities.entry(activity.clone()).or_insert(0) += 1;
            }

            // Track end activity (last in trace)
            if i == activities.len() - 1 {
                *end_activities.entry(activity.clone()).or_insert(0) += 1;
            }

            // Branchless edge collection: iterate over consecutive pairs
            if i + 1 < activities.len() {
                let from_id = node_id;
                let to_activity = &activities[i + 1];
                let to_id = *node_map.entry(to_activity.clone()).or_insert_with(|| {
                    let id = dfg.nodes.len();
                    dfg.nodes.push(DFGNode::new(to_activity.clone(), 0));
                    id
                });

                // Increment edge count (no branching, just or_insert + dereference)
                *edge_counts.entry((from_id, to_id)).or_insert(0) += 1;
            }
        }
    }

    // Materialize edges from integer-keyed counts (single loop, no lookups)
    for ((from_id, to_id), frequency) in edge_counts.iter() {
        let from = &dfg.nodes[*from_id].activity;
        let to = &dfg.nodes[*to_id].activity;
        dfg.edges
            .push(DFGEdge::new(from.clone(), to.clone(), *frequency));
    }

    // Set start and end activities
    dfg.start_activities = start_activities.keys().cloned().collect();
    dfg.end_activities = end_activities.keys().cloned().collect();

    Ok(dfg)
}

#[cfg(test)]
mod tests {
    use super::*;

    #[test]
    fn test_dfg_discovery() {
        let log = EventLog::new(
            vec![Trace::new(
                "case1".to_string(),
                vec![Event::with_activity("A"), Event::with_activity("B")],
            )],
            Vec::new(),
        );

        let dfg = discover_dfg(&log, "concept:name").unwrap();

\end{lstlisting}

\section{PM4PY-LSP Protocol Implementation}
The complete semantic tracking and LSP trait implementation mapping Python typestates to WASM capabilities.
\begin{lstlisting}[language=Rust]
#![doc = "PM4Py Living LSP bridge using tower-lsp-max."]
#![doc = ""]
#![doc = "This crate provides the Language Server Protocol (LSP) implementation"]
#![doc = "for bridging Python's PM4Py library with the wasm4pm ecosystem."]
#![forbid(unsafe_code)]
#![warn(clippy::all)]

use async_trait::async_trait;
use regex::Regex;
use serde::{Deserialize, Serialize};
use serde_json::Value;
use std::collections::HashMap;
use std::path::Path;
use std::sync::{Arc, OnceLock};
use std::time::Duration;
use tokio::fs;
use tokio::sync::Mutex;
use tower_lsp_max::jsonrpc::{Error, Result};
use tower_lsp_max::lsp_types::*;
use tower_lsp_max::max_protocol;
use tower_lsp_max::{Client, LanguageServer};
use uuid::Uuid;
use wasm4pm_compat::hash;

static RE_PARITY_CSV: OnceLock<regex::Regex> = OnceLock::new();

pub mod analysis;
pub mod diagnostics;
pub mod fixtures;
pub mod parity;
pub mod pm4py_bridge;
pub mod receipts;

#[derive(Debug, Serialize, Deserialize, Clone)]
pub struct ParityFixture {
    pub csv_path: String,
    pub parameters: HashMap<String, String>,
    pub expected_outcome: String,
}

#[derive(Debug)]
pub struct Backend {
    pub client: Client,
    pub documents: Arc<Mutex<HashMap<Url, String>>>,
    pub receipts: Arc<Mutex<HashMap<String, max_protocol::Receipt>>>,
    pub pending_scans: Arc<Mutex<HashMap<Url, tokio::task::AbortHandle>>>,
}

impl Backend {
    pub fn new(client: Client) -> Self {
        Self {
            client,
            documents: Arc::new(Mutex::new(HashMap::new())),
            receipts: Arc::new(Mutex::new(HashMap::new())),
            pending_scans: Arc::new(Mutex::new(HashMap::new())),
        }
    }

    async fn on_change(&self, uri: Url, text: String) {
        // Early-return identity check: skip scan if content is unchanged
        {
            let current = self.documents.lock().await.get(&uri).cloned();
            if current.as_deref() == Some(text.as_str()) {
                return; // identical content, skip full scan
            }
        }

        self.documents
            .lock()
            .await
            .insert(uri.clone(), text.clone());

        // Debounce: abort any pending scan for this URI, then schedule a new one
        {
            let mut pending = self.pending_scans.lock().await;
            if let Some(handle) = pending.remove(&uri) {
                handle.abort();
            }
        }

        let pending_scans = self.pending_scans.clone();
        let client = self.client.clone();
        let uri_clone = uri.clone();
        let text_clone = text.clone();

        let handle = tokio::spawn(async move {
            tokio::time::sleep(Duration::from_millis(150)).await;
            pending_scans.lock().await.remove(&uri_clone);
            let diagnostics = diagnose_text(&text_clone);
            client
                .publish_diagnostics(uri_clone, diagnostics, None)
                .await;
        });

        self.pending_scans
            .lock()
            .await
            .insert(uri, handle.abort_handle());
    }


\end{lstlisting}

\section{Tower-LSP-Max Client Handle}
The fully expanded downstream proxy mapping the LSP 3.18 specification. The presence of \texttt{unimplemented!()} macros marks the honest \texttt{PARTIAL\_ALIVE} checkpoint of this research.
\begin{lstlisting}[language=Rust]
use lsp_types::*;
use serde_json::Value;
use crate::client::ClientError;

/// A handle to interact with the connected Language Server.
/// This acts as the outbound proxy for the downstream client.
#[derive(Debug, Clone)]
pub struct ServerHandle {
    // Inner transport channels (to be implemented via tower-service in the builder)
}

impl ServerHandle {
    // --- Lifecycle ---

    pub async fn initialize(&self, _params: InitializeParams) -> Result<InitializeResult, ClientError> { unimplemented!() }
    pub async fn initialized(&self, _params: InitializedParams) { unimplemented!() }
    pub async fn shutdown(&self) -> Result<(), ClientError> { unimplemented!() }
    pub async fn exit(&self) { unimplemented!() }

    // --- Text Document Synchronization ---

    pub async fn did_open(&self, _params: DidOpenTextDocumentParams) { unimplemented!() }
    pub async fn did_change(&self, _params: DidChangeTextDocumentParams) { unimplemented!() }
    pub async fn did_save(&self, _params: DidSaveTextDocumentParams) { unimplemented!() }
    pub async fn did_close(&self, _params: DidCloseTextDocumentParams) { unimplemented!() }

    // --- Language Features ---

    pub async fn hover(&self, _params: HoverParams) -> Result<Option<Hover>, ClientError> { unimplemented!() }
    pub async fn completion(&self, _params: CompletionParams) -> Result<Option<CompletionResponse>, ClientError> { unimplemented!() }
    pub async fn completion_resolve(&self, _params: CompletionItem) -> Result<CompletionItem, ClientError> { unimplemented!() }
    
    pub async fn signature_help(&self, _params: SignatureHelpParams) -> Result<Option<SignatureHelp>, ClientError> { unimplemented!() }
    
    // In lsp-types 0.94, Declaration, Implementation, and TypeDefinition use GotoDefinitionParams and GotoDefinitionResponse.
    pub async fn goto_definition(&self, _params: GotoDefinitionParams) -> Result<Option<GotoDefinitionResponse>, ClientError> { unimplemented!() }
    pub async fn goto_declaration(&self, _params: GotoDefinitionParams) -> Result<Option<GotoDefinitionResponse>, ClientError> { unimplemented!() }
    pub async fn goto_implementation(&self, _params: GotoDefinitionParams) -> Result<Option<GotoDefinitionResponse>, ClientError> { unimplemented!() }
    pub async fn goto_type_definition(&self, _params: GotoDefinitionParams) -> Result<Option<GotoDefinitionResponse>, ClientError> { unimplemented!() }
    
    pub async fn references(&self, _params: ReferenceParams) -> Result<Option<Vec<Location>>, ClientError> { unimplemented!() }
    pub async fn document_highlight(&self, _params: DocumentHighlightParams) -> Result<Option<Vec<DocumentHighlight>>, ClientError> { unimplemented!() }
    pub async fn document_symbol(&self, _params: DocumentSymbolParams) -> Result<Option<DocumentSymbolResponse>, ClientError> { unimplemented!() }
    
    pub async fn code_action(&self, _params: CodeActionParams) -> Result<Option<CodeActionResponse>, ClientError> { unimplemented!() }
    pub async fn code_action_resolve(&self, _params: CodeAction) -> Result<CodeAction, ClientError> { unimplemented!() }
    
    pub async fn code_lens(&self, _params: CodeLensParams) -> Result<Option<Vec<CodeLens>>, ClientError> { unimplemented!() }
    pub async fn code_lens_resolve(&self, _params: CodeLens) -> Result<CodeLens, ClientError> { unimplemented!() }
    
    pub async fn formatting(&self, _params: DocumentFormattingParams) -> Result<Option<Vec<TextEdit>>, ClientError> { unimplemented!() }
    pub async fn range_formatting(&self, _params: DocumentRangeFormattingParams) -> Result<Option<Vec<TextEdit>>, ClientError> { unimplemented!() }
    pub async fn rename(&self, _params: RenameParams) -> Result<Option<WorkspaceEdit>, ClientError> { unimplemented!() }
    pub async fn prepare_rename(&self, _params: TextDocumentPositionParams) -> Result<Option<PrepareRenameResponse>, ClientError> { unimplemented!() }
    
    pub async fn semantic_tokens_full(&self, _params: SemanticTokensParams) -> Result<Option<SemanticTokensResult>, ClientError> { unimplemented!() }
    pub async fn semantic_tokens_full_delta(&self, _params: SemanticTokensDeltaParams) -> Result<Option<SemanticTokensFullDeltaResult>, ClientError> { unimplemented!() }
    pub async fn semantic_tokens_range(&self, _params: SemanticTokensRangeParams) -> Result<Option<SemanticTokensRangeResult>, ClientError> { unimplemented!() }
    
    pub async fn inlay_hint(&self, _params: InlayHintParams) -> Result<Option<Vec<InlayHint>>, ClientError> { unimplemented!() }
    pub async fn inlay_hint_resolve(&self, _params: InlayHint) -> Result<InlayHint, ClientError> { unimplemented!() }

    // --- Workspace Features ---

    pub async fn symbol(&self, _params: WorkspaceSymbolParams) -> Result<Option<Vec<SymbolInformation>>, ClientError> { unimplemented!() }
    pub async fn execute_command(&self, _params: ExecuteCommandParams) -> Result<Option<Value>, ClientError> { unimplemented!() }
    pub async fn did_change_configuration(&self, _params: DidChangeConfigurationParams) { unimplemented!() }
    pub async fn did_change_watched_files(&self, _params: DidChangeWatchedFilesParams) { unimplemented!() }
    pub async fn did_change_workspace_folders(&self, _params: DidChangeWorkspaceFoldersParams) { unimplemented!() }
}

\end{lstlisting}
